\section{Proofs for Section~\ref{sec:fields-certificates}}
\label{app:fields-certificates}

For a rational symmetric matrix $M$ we use the bound
$\norm M\le\max_i\sum_j|M_{ij}|\le\sum_{i,j}|M_{ij}|$, and for $M\succ0$ of
order $h$ the rational margin $\mu_M=\det M/(\tr M)^{h-1}\le\lambda_{\min}(M)$
(Lemma~\ref{lem:models-det-trace}). A positive definite rational Gram matrix on
a vector of rational polynomials yields a sum of squares of rational
polynomials of polynomial size (Lemma~\ref{lem:models-rational-squares}).

\subsection{Proof of Theorem~\ref{thm:fields-multiplier}}
\label{app:fields-multiplier}

Put $X_0=1$ and let $w=(X_jg_i)_{0\le j\le n,\,1\le i\le m}$, a list of
$m(n+1)$ polynomials of degree at most three, possibly with repetitions. Then
$w^{\mathsf T}w=\omega F_0$. Let $\gamma$ be the coordinate vector with
$\gamma^{\mathsf T}w=G$ (entries $\kappa_i$ at the positions $(0,i)$), let
$a_a$ be the coordinate vector with $a_a^{\mathsf T}w=R_a$ given by the
affine multipliers, and let $U$ be the $m(n+1)\times n$ matrix whose $j$th
column represents $X_jG$ (entries $\kappa_i$ at the positions $(j,i)$).
Write $R_a=X^{\mathsf T}T_aX+r_a^{\mathsf T}X+\sigma_a$ with $T_a$ symmetric,
put $b_a=XR_a\in\Q[X]^n$, and let $v=(w,b_1,\ldots,b_s)$. Define
\[
 A_0=c_G\sum_aa_aa_a^{\mathsf T}
   -\tfrac12\sum_a\sigma_a(\gamma a_a^{\mathsf T}+a_a\gamma^{\mathsf T}),
 \qquad
 C_a=\tfrac12\bigl(a_a\ell_G^{\mathsf T}-UT_a-\gamma r_a^{\mathsf T}\bigr),
\]
$C=(C_1\ \cdots\ C_s)$, and $\mathcal H=I_s\otimes H$. Then
\begin{equation}
 v^{\mathsf T}\begin{pmatrix}A_0&C\\C^{\mathsf T}&\mathcal H\end{pmatrix}v=0 .
 \label{eq:fields-zero-identity}
\end{equation}
Indeed, the contribution of target $a$ is
\[
\begin{aligned}
 &c_GR_a^2-\sigma_aGR_a+R_a\,\ell_G^{\mathsf T}XR_a
   -(XG)^{\mathsf T}T_aXR_a-G\,r_a^{\mathsf T}XR_a+R_a^2X^{\mathsf T}HX\\
 &\qquad=R_a^2G-GR_a\bigl(X^{\mathsf T}T_aX+r_a^{\mathsf T}X+\sigma_a\bigr)=0,
\end{aligned}
\]
using $U^{\mathsf T}w=XG$. The terms with $\sigma_a$ and $r_a$ are needed
because the targets need not be homogeneous.

Let $\eta=\mu_H$, $B_A=\sum_{i,j}|(A_0)_{ij}|$, and
\[
 \epsilon=\min\Bigl\{1,\ \frac1{2(1+B_A)},\
   \frac\eta{4(1+\norm C_F^2)}\Bigr\},
 \qquad
 Q_\epsilon=\begin{pmatrix}I+\epsilon A_0&\epsilon C\\
   \epsilon C^{\mathsf T}&\epsilon\mathcal H\end{pmatrix}.
\]
By \eqref{eq:fields-zero-identity}, $v^{\mathsf T}Q_\epsilon v=\omega F_0$. The
top block satisfies $I+\epsilon A_0\succeq\frac12I$, so its inverse has norm
at most two, and the Schur complement is at least
$(\epsilon\eta-2\epsilon^2\norm C_F^2)I\succeq\frac{\epsilon\eta}2I$. Hence
$Q_\epsilon\succ0$, whatever the redundancies in $v$.

Let $A_R=(a_1\ \cdots\ a_s)$ and
$D_J=\diag(A_RJA_R^{\mathsf T},\,J\otimes I_n)$. Then
$v^{\mathsf T}D_Jv=R^{\mathsf T}JR+\norm X^2R^{\mathsf T}JR=\omega R^{\mathsf T}JR$.
With $B_J=\max_i\sum_j|(D_J)_{ij}|$ put
$\lambda_*=\lceil(1+B_J)/\mu_{Q_\epsilon}\rceil$. For $\lambda\ge\lambda_*$,
\[
 \omega(\lambda F_0-R^{\mathsf T}JR)=v^{\mathsf T}(\lambda Q_\epsilon-D_J)v,
 \qquad \lambda Q_\epsilon-D_J\succeq I .
\]
If a rational positive definite full Hessian Gram $M_F$ of $F_0$ is supplied,
let $B_R$ be a rational symmetric full Hessian Gram of $R^{\mathsf T}JR$,
obtained by distributing each coefficient of its Hessian biform over the
pairs of entries of $w(X,v)$ whose product is the corresponding monomial, and
enlarge $\lambda_*$ to at least $(1+\sum_{i,j}|(B_R)_{ij}|)/\mu_{M_F}$; then
$\lambda M_F-B_R\succeq I$. If the $g_i$ vanish at $p$, so do the $R_a$, and
$\lambda F_0-R^{\mathsf T}JR$ is a combination of products of two polynomials
vanishing at $p$, so its value and gradient vanish there.

All matrices have dimension at most $m(n+1)+sn$ and rational entries obtained
from the input by a bounded number of arithmetic operations; $\mu_H$ and
$\mu_{Q_\epsilon}$ are determinants and trace powers of such matrices. Hence
$\lambda_*$ and the Gram certificate have polynomial size, and rational
elimination with the binary splitting of positive weights produces the
squares in polynomial time. A prescribed $\lambda$ contributes its own length.

\subsection{Affine membership is a restriction}
\label{app:fields-membership-example}

For $g_1=x^2+y^2+z^2-1$, $g_2=x^2$, $g_3=xy-z$, $g_4=yz-1$, direct expansion
confirms
\[
 xg_1+(-x+y^2-1)g_2+(-xy+xz-y)g_3+(-x^2-x-1)g_4=1 .
\]
Define the linear functional on polynomials of degree at most three
\[
 \psi(P)=[1]P+[x]P+[yz]P+[z^2]P+[xy^2]P+[xyz]P .
\]
It vanishes on all sixteen products $mg_i$, $m\in\{1,x,y,z\}$: for instance
$\psi(g_1)=-1+1$, $\psi(xg_1)=\psi(x^3+xy^2+xz^2-x)=1-1$,
$\psi(yg_3)=\psi(xy^2-yz)=1-1$, $\psi(zg_3)=\psi(xyz-z^2)=1-1$,
$\psi(g_4)=1-1$, $\psi(xg_4)=\psi(xyz-x)=1-1$, and the remaining products
contain no monomial seen by $\psi$. Since $\psi(1)=1$, the constant $1$ is not
in the span of the $g_i$ and $X_jg_i$. The four quadratics have no common
complex zero: $x^2=0$ forces $x=0$, then $z=0$, contradicting $yz=1$.

\subsection{Proof of Theorem~\ref{thm:fields-higher-membership}}
\label{app:fields-higher}

\paragraph{The lists.}
For $\alpha\in\Z_{\ge0}^n$ with $|\alpha|\le d$ let
$\kappa_\alpha=d!/((d-|\alpha|)!\,\alpha_1!\cdots\alpha_n!)$, a positive
integer, so that $\omega^d=\sum_{|\alpha|\le d}\kappa_\alpha X^{2\alpha}$.
Let $w=(X^\alpha g_i)_{|\alpha|\le d,\,1\le i\le m}$, of level zero, and
$D_0=\diag(\kappa_\alpha)$; then $w^{\mathsf T}D_0w=\omega^dF_0$ and
$D_0\succeq I$. For each target $a$ and each $\alpha$ with
$|\alpha|\le d-1$ add the block $b_{a,\alpha}=X^\alpha XR_a$ of $n$
polynomials, of level $|\alpha|+1$, and let $v$ be the concatenation of $w$
and all blocks. Its length is
$t=m\binom{n+d}d+sn\binom{n+d-1}{d-1}$.

Every \emph{parent} $X^\alpha R_a$, $|\alpha|\le d-1$, has an explicit
representation: for $\alpha=0$ by the membership coordinates in $w$, and for
$|\alpha|=e\ge1$ as the $j$th entry of $b_{a,\alpha-e_j}$, of level $e$, for
any $j$ with $\alpha_j>0$. The polynomials $X^\alpha G$ and $X^\alpha XG$ are
represented in $w$, because $G$ is a constant combination of the $g_i$ and
$|\alpha|+1\le d$.

\paragraph{Zero identities.}
Write $R_a=X^{\mathsf T}T_aX+r_a^{\mathsf T}X+\sigma_a$. For fixed
$(a,\alpha)$ put $u=X^\alpha$, $\pi=uR_a$, $\chi=uG$, $\zeta=uXR_a$, and
$\Upsilon=uXG$. Then
\begin{equation}
 \zeta^{\mathsf T}H\zeta+\pi\,\ell_G^{\mathsf T}\zeta+c_G\pi^2
 -\Upsilon^{\mathsf T}T_a\zeta-\chi\,r_a^{\mathsf T}\zeta-\sigma_a\chi\pi=0,
 \label{eq:fields-level-identity}
\end{equation}
because the first three terms equal $u^2GR_a^2$ and the last three equal
$-u^2GR_a^2$. Let $Z_{a,\alpha}$ be the symmetric rational matrix on $v$ that
represents the left side through the representations above, cross terms split
equally. Its block on $b_{a,\alpha}=\zeta$ is $H$; let $E_{a,\alpha}$ be
$Z_{a,\alpha}$ with that block removed. With $e=|\alpha|$, the entries of
$E_{a,\alpha}$ join the following levels: $c_G\pi^2$ joins $(e,e)$;
$\sigma_a\chi\pi$ joins $(0,e)$; $\pi\ell_G^{\mathsf T}\zeta$ joins
$(e,e+1)$; and the two terms with $\Upsilon$ and $\chi$ join $(0,e+1)$.

\paragraph{One perturbation.}
Let $K=\sum_{a,\alpha}\sum_{i,j}|(E_{a,\alpha})_{ij}|$,
$\eta_0=\min\{1,\mu_H\}$, and $\rho=\eta_0/(2(1+K))\le\frac12$. Put
\[
 Q=(D_0\oplus0)+\sum_{a,\alpha}\rho^{2(|\alpha|+1)}Z_{a,\alpha},
\]
so $v^{\mathsf T}Qv=\omega^dF_0$. Let $S$ be diagonal with entry $\rho^{-l}$
at every coordinate of level $l$. In $SQS$ the blocks $D_0$ and one copy of
$H$ on every new block appear unchanged, because
$\rho^{2(e+1)}\rho^{-2(e+1)}=1$. An entry of $E_{a,\alpha}$ joining levels
$l,l'$ is multiplied by $\rho^{2(e+1)-l-l'}$, and the four level pairs give
exponents $2$, $e+2$, $1$, and $e+1$, all at least one. Hence the corrections
in $SQS$ have total absolute entry sum at most $K\rho<\eta_0/2$, and
$SQS\succeq\eta_0I-\frac{\eta_0}2I$. Since levels are at most $d$,
\[
 Q\succeq\tfrac12\eta_0\rho^{2d}I .
\]

\paragraph{The target.}
For $|\beta|\le d$ let $V_\beta$ be the rational $s\times t$ matrix expressing
the vector $X^\beta R$ in $v$, through the membership coordinates for
$\beta=0$ and through an entry of $b_{a,\beta-e_j}$ otherwise. Then
$D_J=\sum_{|\beta|\le d}\kappa_\beta V_\beta^{\mathsf T}JV_\beta$ represents
$\omega^dR^{\mathsf T}JR$. With $B_J=\max_i\sum_j|(D_J)_{ij}|$ and
$\lambda_*=\lceil2(B_J+1)/(\eta_0\rho^{2d})\rceil$, every rational
$\lambda\ge\lambda_*$ gives
$\omega^d(\lambda F_0-R^{\mathsf T}JR)=v^{\mathsf T}(\lambda Q-D_J)v$ with
$\lambda Q-D_J\succeq I$. The entries of $v$ have degree at most $d+2$.
Because $v$ contains every $g_i$, this identity gives
$\omega^d(\lambda F_0-R^{\mathsf T}JR)\ge\sum_ig_i^2$, so a real zero of
$\lambda F_0-R^{\mathsf T}JR$ is a common zero of the $g_i$; conversely the
membership representations make every $R_a$ vanish there.

\paragraph{Size and denominators.}
The matrices have dimension $t$, polynomial in $\binom{n+d}d$, $n$, $m$, $s$;
the multinomials have $O(d\log(n+1))$ bits; $K$ is a sum of polynomially many
explicit rationals; and $\rho^{2d}$ has polynomial length. This gives the
complexity claim, with no successive determinant bounds. If $d$ is even, the
identity gives a rational-function sum of squares with denominator
$\omega^{d/2}$; if $d$ is odd, multiply once more by $\omega$ as in
\eqref{eq:fields-multiplier-to-denominator} and use $\omega^{(d+1)/2}$. For
$d=0$ the targets lie in the constant span of the $g_i$, and direct domination
of their Gram matrix on $(g_i)$ gives a polynomial certificate.

\subsection{The explicit multiplier certificate of
Example~\ref{ex:fields-ternary-multiplier}}
\label{app:fields-ternary-multiplier}

Let
\[
\begin{aligned}
 b=(&z^2-\tfrac x2,\ y^2-2z,\ xz-\tfrac y2,\ xy-1,\ x^2-yz,\
     z^3-\tfrac y4,\ yz^2-\tfrac12,\ y^2z-x,\\
    &y^3-2yz,\ xz^2-\tfrac{yz}2,\ xyz-z,\ xy^2-y,\ x^2z-\tfrac12,\
     x^2y-x,\ x^3-z)^{\mathsf T}
\end{aligned}
\]
and let $N$ be the symmetric integer matrix whose rows are
{\scriptsize
\[
\setlength{\arraycolsep}{2pt}
\begin{array}{rrrrrrrrrrrrrrr}
4240&527&792&-124&992&-1728&696&-1112&216&-312&1212&-384&-1200&76&-8\\
527&556&376&-194&361&-16&136&-92&0&280&-340&-232&44&168&-128\\
792&376&3176&-175&568&-1000&584&-232&256&-1272&432&-912&-1112&990&-384\\
-124&-194&-175&980&-72&432&-1080&552&-200&104&-408&488&-320&-624&136\\
992&361&568&-72&1960&-216&480&400&-80&-912&184&-400&-576&440&-720\\
-1728&-16&-1000&432&-216&2624&-1440&700&-24&0&-912&8&1052&-256&-40\\
696&136&584&-1080&480&-1440&3784&-1896&440&-240&728&-680&-1344&1232&-352\\
-1112&-92&-232&552&400&700&-1896&2000&-480&-96&-832&400&560&-816&240\\
216&0&256&-200&-80&-24&440&-480&320&-24&240&-384&56&264&-80\\
-312&280&-1272&104&-912&0&-240&-96&-24&3368&-672&264&40&-488&460\\
1212&-340&432&-408&184&-912&728&-832&240&-672&2800&-552&-952&376&-304\\
-384&-232&-912&488&-400&8&-680&400&-384&264&-552&1296&24&-944&244\\
-1200&44&-1112&-320&-576&1052&-1344&560&56&40&-952&24&2760&-528&0\\
76&168&990&-624&440&-256&1232&-816&264&-488&376&-944&-528&1848&-640\\
-8&-128&-384&136&-720&-40&-352&240&-80&460&-304&244&0&-640&832
\end{array}
\]}
Expanding both sides shows $b^{\mathsf T}Nb=8(1+x^2+y^2+z^2)F_\star$. The
leading principal minors of $N-8I$ of orders $1,\ldots,15$ are
{\small
\[
\begin{gathered}
4232,\quad 2041407,\quad 5839006240,\quad 5263436936513,\quad
8309698919905512,\\
11480475166599386752,\quad 24140997007816543732224,\quad
9072065142124268516803584,\\
1221674797258437785568657408,\quad
1569452470480427736351156994048,\\
1967046692952416690050593546305536,\quad
844438887950283618215158211684597760,\\
196739716166438745148196850208516079616,\quad
21440610749731697832005735132127356256256,\\
497197930466744687768853864373192977547264,
\end{gathered}
\]}
all positive, so $N-8I\succ0$ by Sylvester's criterion and $N/8$ is a rational
positive definite Gram matrix of $\omega F_\star$ on $b$. Each entry of $b$
vanishes at $p_\star$, by the evaluation rule
$x^iy^jz^h\mapsto2^{\lfloor e/5\rfloor-i-h}a^{e\bmod5}$ with $e=4i+j+2h$;
the fifteen entries have distinct leading monomials, none of which occurs in
another entry, so they are independent; and the rational polynomials of
degree at most three vanishing at $p_\star$ form a space of dimension
$20-5=15$. Thus $b$ is a basis of that space, which explains why a certificate
of this shape can exist. The identity and the positivity are the only finite
verifications in Example~\ref{ex:fields-ternary-multiplier}.

\subsection{Proofs of Theorems~\ref{thm:fields-radial}
and~\ref{thm:fields-radial-height}}
\label{app:fields-radial}

Throughout, $p=(a^2,a^3,a^4)$ with $a=2^{1/5}$, the variables are
$X=(x_0,x_1,x_2)$, and $f^\circ$ is as in Section~\ref{sec:fields-radial}.

\paragraph{Theorem~\ref{thm:fields-radial}(a) and (b).}
Let $A^\circ$ be a rational positive definite full Hessian Gram of
$f^\circ$. Since $\nabla^2f_t(X)=\nabla^2f^\circ(tX)$ and
$w(tX,v)=\diag(I_3,tI_9)\,w(X,v)$, the matrix
$\diag(I_3,tI_9)A^\circ\diag(I_3,tI_9)$ is a rational positive definite full
Hessian Gram of $f_t$, and $\nabla^2f_t\succeq I$. The minimizer is $p/t$, with
$f_t(p/t)=0$ and $\nabla^2f_t(p/t)=\nabla^2f^\circ(p)$. Each coefficient of
$f_t$ is $t^{|\gamma|-2}[X^\gamma]f^\circ$ with $|\gamma|\le4$, so
$L(t)=O(\log t)$; the quartic part of $f^\circ$ is nonzero, so some
coefficient is a nonzero integer multiple of $t^2$ and $L(t)\ge\log_2t$. For
(b), write $\omega f^\circ=\sum_jg_j^2$ with fixed rational cubics. Then
\[
 (1+t^2\norm X^2)f_t(X)=t^{-2}\omega(tX)f^\circ(tX)
 =\sum_j\bigl(g_j(tX)/t\bigr)^2,
\]
a fixed number of rational cubics with coefficients of $O(\log t)$ bits.
The identity \eqref{eq:fields-multiplier-to-denominator}, with $\omega$
replaced by $1+t^2\norm X^2$, gives the rational-function certificate. If
$f_t$ were a rational sum of squares, the substitution $X=x/t$ would make
$f^\circ$ one, contradicting Theorem~\ref{thm:fields-prime} for $f^\circ$;
Lemma~\ref{lem:fields-affine-denominator} then excludes degree at most one.

\paragraph{Vanishing spaces.}
For $d\ge2$ let $I_d=\mathcal I_d^\Q(p)$ and let $V_d$ be its real span. Use
the representatives $\varrho_0=1$, $\varrho_1=x_0x_2$, $\varrho_2=x_0$,
$\varrho_3=x_1$, $\varrho_4=x_2$, whose values at $p$ are
$1,2a,a^2,a^3,a^4$. A monomial $m=x_0^ix_1^jx_2^h$ has
$m(p)=a^e=2^{\lfloor e/5\rfloor}a^{e\bmod5}$ with $e=2i+3j+4h$, so
$m(p)=c_m\varrho_{e\bmod5}(p)$ with $c_m=2^{\lfloor e/5\rfloor-[e\equiv1]}$.
For every monomial $m$ of degree at most $d$ other than the representatives,
put $\pi_m=m-c_m\varrho_{e\bmod5}$. These polynomials form a basis of $I_d$:
each contains its own monomial $m$ with coefficient one and no other
non-representative monomial, and a rational combination of the
representatives vanishing at $p$ is zero because $1,a,\ldots,a^4$ are
linearly independent. The bases are nested in $d$. For $\deg m\le d$ we have
$e\le4d$, so $c_m\in\{2^{-1},1,\ldots,2^{d-1}\}$; every $\pi_m$ has
coefficients in $\frac12\Z$ and $\ell_1$-norm at most $2^d$. For $d\ge3$ the
new basis elements, those with $\deg m=d$, have leading form exactly $m$.
Finally, $\varphi_5(g^2)\ge0$ for every $g\in V_2$: every rational element of
$I_2$ satisfies the linear condition $v_2+h_{00}=0$ of
Lemma~\ref{lem:fields-prime-functional}, hence so does every real combination,
and then $\varphi_5(g^2)=v_0^2+2g_0(v_2+h_{00})=v_0^2$.

\paragraph{Grid functionals.}
A real polynomial in three variables of degree at most $d$ in each variable
that vanishes on $\{0,1,\ldots,d\}^3$ is zero: for fixed values of the last
two variables in the grid the univariate polynomial in the first variable has
$d+1$ roots, so its coefficients vanish on the grid in the remaining variables,
and induction finishes. Let $E_2(P)=\sum_{u\in\{0,1,2\}^3}P(u)$ and, for
$d\ge3$, let $M_d(P)=\sum_{u\in\{0,\ldots,d\}^3}P_{2d}(u)$, where $P_{2d}$ is
the homogeneous part of degree $2d$.

\paragraph{Separators.}
Let $K_0=E_2(f^\circ)+1$, a positive integer since $f^\circ\ge0$ has integer
coefficients, and
$\psi_2=\varphi_5+E_2/K_0$ on $\R[X]_{\le4}$. Then
$a_0=-\psi_2(f^\circ)=4-E_2(f^\circ)/K_0>3$. For $0\ne g\in V_2$, $g$ has
degree at most two in each variable, so $E_2(g^2)=\sum_ug(u)^2>0$, and
$\psi_2(g^2)>0$. All monomial values of $\psi_2$ lie in $K_0^{-1}\Z$.

Let $d\ge3$ and suppose $\psi_{d-1}$ is a linear functional on
$\R[X]_{\le2d-2}$ with monomial values in $K_0^{-1}\Z$, positive on squares
of nonzero elements of $V_{d-1}$, and with $\psi_{d-1}(f^\circ)=-a_0$. Extend
it by zero to the monomials of degrees $2d-1$ and $2d$, obtaining
$\hat\psi_{d-1}$. On the nested basis of $I_d$ its Gram matrix is
$\left(\begin{smallmatrix}A&B\\B^{\mathsf T}&C\end{smallmatrix}\right)$, where
$A\succ0$ is the Gram matrix of $\psi_{d-1}$ on the basis of $I_{d-1}$. The
functional $M_d$ vanishes on products of an old basis element with any basis
element, which have degree at most $2d-1$; on two new elements with leading
monomials $m,m'$ it gives
$D_{mm'}=\sum_um(u)m'(u)$. The integer matrix $D$ is positive definite: for a
nonzero coefficient vector the form $\sum c_mm$ is nonzero, has degree at most
$d$ in each variable, and does not vanish on the grid. Put
$S=C-B^{\mathsf T}A^{-1}B$, $m_d=\max_i\sum_j|S_{ij}|$,
$\beta_d=\mu_D$, $T_d=\lceil(m_d+1)/\beta_d\rceil$, and
$\psi_d=\hat\psi_{d-1}+T_dM_d$. The Gram matrix of $\psi_d$ on the basis of
$I_d$ is $\left(\begin{smallmatrix}A&B\\B^{\mathsf T}&C+T_dD\end{smallmatrix}\right)$,
with Schur complement $S+T_dD\succeq I$, so it is positive definite. Its
monomial values remain in $K_0^{-1}\Z$, and $\psi_d(f^\circ)=-a_0$ because
$\deg f^\circ=4\le2d-2$.

\paragraph{Theorem~\ref{thm:fields-radial}(c).}
Fix $N\ge0$, put $d=N+2$, $s=\norm X^2$, and
$A_N=\sum_{j=1}^N\binom Nj|\psi_d(s^jf^\circ)|$, and let
$\tau_N=\bigl\lceil\sqrt{\max\{1,2(A_N+1)/a_0\}}\,\bigr\rceil$. For an
integer $t\ge\tau_N$ put $\epsilon=t^{-2}\le\min\{1,a_0/(2(A_N+1))\}$. Then
\[
 \psi_d\bigl((1+\epsilon s)^Nf^\circ\bigr)\le-a_0+\epsilon A_N<-a_0/2 .
\]
If $\omega^Nf_t$ were a sum of squares in $\Q[X]$, the substitution $X=x/t$ and
multiplication by $t^2$ would make $(1+\epsilon s)^Nf^\circ$ one. Its factors
have degree at most $d$ and vanish at $p$, so they lie in $I_d$, and
$\psi_d$ would be nonnegative on it. This contradiction proves the first
claim; $\tau_N$ is computed by exact rational arithmetic. For a fixed
$u\in\Q[X]$ with $u(0)>0$, let $d=2+\lceil\deg u/2\rceil$. As $t\to\infty$
the polynomial $u(x/t)f^\circ(x)$ converges coefficientwise to
$u(0)f^\circ$, so $\psi_d(u(x/t)f^\circ)\to-u(0)a_0<0$, and the same argument
excludes a rational sum of squares of $uf_t$ for large $t$.

\paragraph{Theorem~\ref{thm:fields-radial-height}.}
Let $H_d=\max\{1,\max_{\deg m\le2d}|K_0\psi_d(m)|\}$, an integer. Fix
$d\ge3$, and let $n_d=\dim I_{d-1}\le d^3$ and $h_d=\binom{d+2}2\le2d^2$ be
the numbers of old and new basis elements. A product of two basis elements
has coefficients in $\frac14\Z$ and $\ell_1$-norm at most $2^{2d}$, so every
entry of the Gram matrix of $\hat\psi_{d-1}$ becomes an integer of absolute
value at most $U_d=2^{2d+2}H_{d-1}$ after multiplication by $Q_0=4K_0$. The
integer matrix $Q_0A$ is positive definite, so $\det(Q_0A)\ge1$, and its
cofactors are at most $(n_d-1)!\,U_d^{n_d-1}$. Hence
$|(A^{-1})_{ij}|\le Q_0(n_dU_d)^{n_d-1}$, and
\[
 |S_{ij}|\le\frac{U_d}{Q_0}+n_d^2\Bigl(\frac{U_d}{Q_0}\Bigr)^2Q_0(n_dU_d)^{n_d-1}
 \le2(n_dU_d)^{n_d+1},\qquad m_d\le2h_d(n_dU_d)^{n_d+1}.
\]
The entries of $D$ are at most $J_d=(d+1)^3d^{2d}$, $\det D\ge1$, and
$\tr D\le h_dJ_d$, so $\beta_d^{-1}\le(h_dJ_d)^{h_d-1}$ and
$T_d\le(m_d+2)(h_dJ_d)^{h_d}$. The new monomial values satisfy
$|K_0T_dM_d(m)|\le K_0T_dJ_d$, and old values are unchanged, so
$H_d\le\max\{H_{d-1},K_0(m_d+2)(h_dJ_d)^{h_d+1}\}$. Taking logarithms, with
$\log_2U_d=2d+2+\log_2H_{d-1}$, $n_d+1\le2d^3$, and $\log_2J_d\le5d\log_2(d+1)$,
there is an absolute constant $C_1$ with
\[
 \log_2H_d\le C_1d^3\bigl(d+\log_2H_{d-1}\bigr)+\log_2K_0 .
\]
Put $\beta'_d=d+\log_2H_d+\log_2K_0\ge1$. Then
$\beta'_d\le C_2d^3\beta'_{d-1}$ for an absolute constant $C_2$, so
$\beta'_d\le C_2^d(d!)^3\beta'_2$ and $\log_2\beta'_d\le C_3d\log_2d$ for
$d\ge3$, with $C_3$ depending only on $f^\circ$.

For the threshold, $|\psi_d(P)|\le(H_d/K_0)\norm P_1$, where $\norm P_1$ is
the sum of the absolute coefficients, and $\norm{s^jf^\circ}_1\le
3^j\norm{f^\circ}_1$. Hence $A_N\le6^NH_{N+2}\norm{f^\circ}_1$, and since
$a_0>3$,
$\log_2\tau_N\le C_4\beta'_{N+2}$ and
$\log_2\log_2\tau_N\le C_5(N+2)\log_2(N+2)$ for $N\ge1$, with constants
depending only on $f^\circ$. The bound also holds for
$\tau'_N=\max_{j\le N}\tau_j$, since its right side increases with $N$.

For an integer $t$ put $u=\log_2\log_2t$, and assume $u$ is large enough
that $N=\lfloor u/(2C_5\log_2u)\rfloor-2\ge1$. Then $N+2\le u$ and
$C_5(N+2)\log_2(N+2)\le u/2$, so $t\ge\tau'_N$. By part (c), $\omega^jf_t$ is
not a rational sum of squares for any $j\le N$, so $\nu(f_t)\ge N+1$. This
gives $\nu(f_t)\ge c\,u/\log_2u$ for $t\ge t_0$, which is the claim, and
$L(t)=\Theta(\log t)$ gives the form in $L$. The argument does not use the
finiteness of $\nu(f_t)$, which is Proposition~\ref{prop:fields-radial-finite}.

\subsection{Proof of Proposition~\ref{prop:fields-radial-finite}}
\label{app:fields-radial-finite}

\paragraph{Imported results.}
All rings are commutative and contain $\Q$. A \emph{cone} in an additive group
is closed under addition and contains zero. For a cone $C$ in a group $G$, an
element $u\in C$ is an \emph{order unit} of $(G,C)$ if for every $g\in G$ some
positive integer $r$ has $ru\pm g\in C$. A \emph{state} of $(G,C,u)$ is an
additive map $\psi:G\to\R$ that is nonnegative on $C$ and has $\psi(u)=1$; a
\emph{pure} state is an extreme point of the convex set of states. A
\emph{preordering} $S$ of a ring $A$ is a cone containing all squares and
closed under multiplication; it is \emph{archimedean} if for every $a\in A$
some positive integer $R$ has $R-a\in S$. For an ideal $I$, an
\emph{$S$-pseudomodule} in $I$ is a cone $C\subseteq I$ with $SC\subseteq C$.
We use the following results of Burgdorf, Scheiderer, and Schweighofer.
\begin{enumerate}[label=(\Roman*)]
\item \cite[Theorem~2.5]{BurgdorfScheidererSchweighofer2012} If $u$ is an
  order unit of $(G,C)$ and $\psi(g)>0$ for every pure state $\psi$ of
  $(G,C,u)$, then $rg\in C$ for some positive integer $r$.
\item \cite[Corollary~4.12]{BurgdorfScheidererSchweighofer2012} Let $S$ be an
  archimedean preordering of $A$, $I$ an ideal, $C\subseteq I$ an
  $S$-pseudomodule with order unit $u$, and $\psi$ a pure state of
  $(I,C,u)$. If $\psi(u^2)\ne0$, there is a ring homomorphism
  $\lambda:A\to\R$, nonnegative on $S$, with $\lambda(u)\ne0$ and
  $\psi(h)=\lambda(h)/\lambda(u)$ for all $h\in I$. If $\psi(u^2)=0$, there
  is a ring homomorphism $\lambda:A\to\R$, nonnegative on $S+I$, with
  $\lambda(I)=0$ and $\psi(ah)=\lambda(a)\psi(h)$ for all $a\in A$ and
  $h\in I$.
\item \cite[Theorem~6.2]{BurgdorfScheidererSchweighofer2012} If $S$ is an
  archimedean preordering of $A$ and $\lambda(h)>0$ for every ring
  homomorphism $\lambda:A\to\R$ nonnegative on $S$, then $rh\in S$ for some
  positive integer $r$.
\end{enumerate}
The order-unit statement of
\cite[Proposition~5.3(a) and Remark~5.5(1)]{BurgdorfScheidererSchweighofer2012}
is also used; we prove the case needed below directly. Every positive
rational is a sum of rational squares, so a preordering of sums of squares is
closed under multiplication by positive rationals, and $rh\in S$ implies
$h\in S$.

\paragraph{The sphere ring.}
Let $y=(y_0,\ldots,y_n)$, $q=\sum_iy_i^2$, $A=\Q[y]/(q-1)$, and let
$S\subseteq A$ be the set of sums of squares, a preordering. Let $B$ be the
set of $a\in A$ with $N\pm a\in S$ for some positive integer $N$. It contains $\Q$,
and it is closed under addition, negation, and multiplication by rationals.
If $N\pm a\in S$, then
\[
 N^2-a^2=\frac1{2N}\bigl((N+a)^2(N-a)+(N-a)^2(N+a)\bigr)\in S
\]
and $N^2+a^2\in S$, so $a^2\in B$; hence $ab=((a+b)^2-(a-b)^2)/4\in B$ for
$a,b\in B$, and $B$ is a subring. In $A$,
$1\pm y_i=\frac12\bigl((1\pm y_i)^2+\sum_{j\ne i}y_j^2\bigr)\in S$, so every
$y_i$ lies in $B$, $B=A$, and $S$ is archimedean. The ring homomorphisms
$A\to\R$ are the evaluations at points of the real unit sphere $S^n$, and all
of them are nonnegative on $S$.

Let $H(y)=y_0^{2e}f(y_1/y_0,\ldots,y_n/y_0)$ be the homogenization of $f$,
a form of degree $2e$. For $y_0=0$ and $y\ne0$, $H(y)$ is the positive
definite leading form of $f$ at $(y_1,\ldots,y_n)$, and for $y_0\ne0$,
$H(y)=y_0^{2e}f(y'/y_0)\ge0$. Hence $H\ge0$ on $\R^{n+1}$, and its zeros on
$S^n$ form the finite set $Z=\{\pm(1,z)/\norm{(1,z)}:f(z)=0\}$. If $Z$ is
empty, (III) gives $H\in S$; continue with the last paragraph.

\paragraph{Algebraicity of the zeros.}
Let $z$ be a real zero of $f$ and $K=\Q(z_1,\ldots,z_n)$. Since $f\ge0$,
$\nabla f(z)=0$, and $\nabla^2f(z)$ is invertible. Suppose $K$ had positive
transcendence degree over $\Q$. Choose a transcendence basis and the
derivation of the rational function field that sends one basis element to
one and the others to zero; in characteristic zero it extends uniquely to a
derivation $\partial$ of the finite algebraic extension $K$. Applying
$\partial$ to the rational identities $\nabla f(z)=0$ gives
$\nabla^2f(z)\,\partial z=0$, hence $\partial z=0$, so $\partial$ vanishes on
the generators of $K$ and is zero, a contradiction. Thus $z$, and every point
of $Z$, has algebraic coordinates.

\paragraph{The vanishing ideal.}
For $x\in Z$ the evaluation $a\mapsto a(x)$ maps $A$ onto the number field
$\Q(x)$, so its kernel $\mathfrak m_x$ is a maximal ideal. Let
$\mathfrak m_1,\ldots,\mathfrak m_t$ be the distinct ideals among them, and
$J=\mathfrak m_1\cap\cdots\cap\mathfrak m_t$. Distinct maximal ideals are
comaximal, and so are their powers: if $a+b=1$ with $a\in\mathfrak m_i$,
$b\in\mathfrak m_j$, expanding $(a+b)^3$ shows
$1\in\mathfrak m_i^2+\mathfrak m_j^2$. Hence
$J=\mathfrak m_1\cdots\mathfrak m_t$ and
$\mathfrak m_1^2\cap\cdots\cap\mathfrak m_t^2
=\mathfrak m_1^2\cdots\mathfrak m_t^2=J^2$. A real point of $S^n$ at which
all of $J$ vanishes gives a ring homomorphism $A/J\cong\prod_iA/\mathfrak m_i\to\R$,
which factors through one residue field; since $H$ lies in every
$\mathfrak m_i$, it vanishes there, so the point is in $Z$. The real zero set
of $J$ on $S^n$ is therefore exactly $Z$.

We show $H\in J^2$. Let $x\in Z$ and $\mathfrak M=\{g\in\Q[y]:g(x)=0\}$, the
preimage of $\mathfrak m_x$. Since $H\ge0$ on $\R^{n+1}$ and $H(x)=0$, the
point $x$ is a minimizer and $\nabla H(x)=0$. Now a polynomial
$g\in\mathfrak M$ with $\nabla g(x)=0$ lies in $\mathfrak M^2$. Indeed, after
an invertible rational linear change of coordinates, which transforms
gradients linearly, the primitive element theorem lets us assume
$\Q(x)=\Q(x_0)$. Let $\mu$ be the minimal polynomial of $x_0$, so
$\mu'(x_0)\ne0$, and write $x_i=c_i(x_0)$ with $c_i\in\Q[T]$. The polynomials
$g_0=\mu(y_0)$ and $g_i=y_i-c_i(y_0)$, $1\le i\le n$, vanish at $x$, and the
quotient of $\Q[y]$ by the ideal they generate is $\Q[y_0]/(\mu)$, a field;
so they generate $\mathfrak M$. Their gradients at $x$, namely
$\mu'(x_0)e_0$ and $e_i-c_i'(x_0)e_0$, are linearly independent over
$\Q(x)$. If $g=\sum_ja_jg_j$ and $\nabla g(x)=\sum_ja_j(x)\nabla g_j(x)=0$,
then every $a_j(x)=0$, so $a_j\in\mathfrak M$ and $g\in\mathfrak M^2$. The
same computation shows that $\mathfrak M/\mathfrak M^2$ has dimension $n+1$
over $\Q(x)$, the class of $g$ corresponding to $\nabla g(x)$. Applied to
$H$, this gives $H\in\mathfrak M^2$, so the class of $H$ lies in
$\mathfrak m_x^2$ for every $x\in Z$, hence in $J^2$.

\paragraph{An order unit.}
Let $b_1,\ldots,b_r$ generate $J$, put $I=J^2$, $C=S\cap I$, and
$u=\sum_ib_i^2$. The cone $C$ is an $S$-pseudomodule in $I$. Every $h\in I$
is a finite sum of terms $ab_ib_j$ with $a\in A$, and for an integer
$R\ge1$ with $R-a^2\in S$,
\[
 R(b_i^2+b_j^2)\pm2ab_ib_j
 =(ab_i\pm b_j)^2+(R-a^2)b_i^2+(R-1)b_j^2\in C .
\]
Summing over a representation of $h$, and using $u-b_i^2\in C$, gives a
positive integer $N$ with $Nu\pm h\in C$. Since also $u\in C$, $u$ is an
order unit of $(I,C)$. The zeros of $u$ on $S^n$
are exactly $Z$.

\paragraph{Pure states.}
Let $\psi$ be a pure state of $(I,C,u)$ and $\lambda$ the homomorphism given
by (II), evaluation at a point $x\in S^n$.

If $\psi(u^2)\ne0$, then $\lambda(u)\ne0$; as $u$ is a sum of squares,
$\lambda(u)>0$, so $x\notin Z$ and $\psi(H)=H(x)/u(x)>0$.

If $\psi(u^2)=0$, then $\lambda(I)=0$. For $b\in J$ we have $b^2\in I$, so
$\lambda(b)^2=0$; thus $\lambda(J)=0$ and $x\in Z$. Let
$\mathfrak m=\mathfrak m_x$ and $K=A/\mathfrak m\cong\Q(x)$, embedded in $\R$
by $\lambda$. The relation $\psi(ah)=\lambda(a)\psi(h)$ shows that $\psi$
vanishes on $\mathfrak mI$ and is $K$-semilinear. The $K$-vector spaces
$I/\mathfrak mI$ and $\mathfrak m^2/\mathfrak m^3$ coincide: both are
annihilated by $\mathfrak m$, hence unchanged by localization at
$\mathfrak m$, and $J$ and $\mathfrak m$ have the same localization because
the other $\mathfrak m_i$ contain elements outside $\mathfrak m$. The space
$\mathfrak m/\mathfrak m^2=\mathfrak M/(\mathfrak M^2+(q-1))$ has dimension
$n$ over $K$, because the class of $q-1$ corresponds to $\nabla q(x)=2x\ne0$;
its elements correspond to the restrictions to $T_xS^n=x^\perp$ of the
$K$-rational covectors $\nabla g(x)$. If $c_1,\ldots,c_n\in\mathfrak m$ have
classes forming a $K$-basis of $\mathfrak m/\mathfrak m^2$, every product of
two elements of $\mathfrak m$ is congruent modulo $\mathfrak m^3$ to a
combination of the $c_ic_j$ with coefficients in $A$, so
$\dim_K\mathfrak m^2/\mathfrak m^3\le n(n+1)/2$.

For $b\in\mathfrak m^2$ let $\mathcal Q_b$ be the Hessian at $x$ of the
restriction of $b$ to $S^n$, a quadratic form on $T_xS^n$, well defined
because $b$ and its tangential differential vanish at $x$. The map
$b\mapsto\mathcal Q_b$ vanishes on $\mathfrak m^3$, satisfies
$\mathcal Q_{kb}=\lambda(k)\mathcal Q_b$, and sends $cc'$ with
$c,c'\in\mathfrak m$ to $dc\otimes dc'+dc'\otimes dc$, where $dc$ is the
tangential differential at $x$; these differentials span the real cotangent
space. So $b\mapsto\mathcal Q_b$ induces a surjective $\R$-linear map from
$(\mathfrak m^2/\mathfrak m^3)\otimes_{K,\lambda}\R$ onto the quadratic forms on
$T_xS^n$, a space of dimension $n(n+1)/2$, at least that of the source; it is
therefore bijective. Hence there is a real symmetric form $\Phi$ on $T_xS^n$
with $\psi(h)=\ip{\Phi}{\mathcal Q_h}$ for all $h\in I$.

The form $\Phi$ is positive semidefinite and nonzero. For $g\in J$ we have
$g^2\in C$, so $0\le\psi(g^2)=2\ip{\Phi}{dg\otimes dg}$. Since $J$ maps onto
$\mathfrak m/\mathfrak m^2$, the differentials $dg$, $g\in J$, fill the
$K$-rational points of the cotangent space, which are dense, so
$\Phi\succeq0$. Also $1=\psi(u)=2\sum_i\ip{\Phi}{db_i\otimes db_i}$, so
$\Phi\ne0$. Finally $\mathcal Q_H\succ0$. Write $x=(x_0,x')$ with $x_0\ne0$
and $z=x'/x_0$, so $f(z)=0$, $\nabla f(z)=0$, and $\nabla^2f(z)\succ0$. With
$\pi(y)=(y_1/y_0,\ldots,y_n/y_0)$ and $H=y_0^{2e}f\circ\pi$ near $x$, the
vanishing value and gradient of $f\circ\pi$ at $x$ give
$\nabla^2H(x)=x_0^{2e}D\pi(x)^{\mathsf T}\nabla^2f(z)D\pi(x)$. The kernel of
$D\pi(x)$ is $\R x$, which meets $x^\perp$ only in zero, and since
$\nabla H(x)=0$ the Hessian of $H|_{S^n}$ at $x$ is the restriction of
$\nabla^2H(x)$ to $x^\perp$. Hence $\mathcal Q_H\succ0$, and
$\psi(H)=\ip{\Phi}{\mathcal Q_H}>0$, a nonzero positive semidefinite form
paired with a positive definite one.

Every pure state is therefore positive at $H$. By (I), $rH\in C\subseteq S$
for some positive integer $r$, and so $H\in S$.

\paragraph{From the sphere to the radial multiplier.}
Write $H=\sum_ip_i^2$ in $A$, with polynomials $p_i\in\Q[y]$, so
$H-\sum_ip_i^2\in(q-1)$. Split $p_i=e_i+o_i$ into its parts that are even and
odd under $y\mapsto-y$. Since $H$ and $q-1$ are even, also
$H-\sum_ip_i(-y)^2\in(q-1)$, and averaging gives
$H-\sum_i(e_i^2+o_i^2)\in(q-1)$. Let $D\ge2e$ be an even integer bounding all
$\deg p_i$. Multiplying each monomial of $e_i$ of degree $2j$ by
$q^{(D-2j)/2}$, and each monomial of $o_i$ of degree $2j+1$ by
$q^{(D-2j-2)/2}$, gives forms $E_i$ of degree $D$ and $O_i$ of degree $D-1$
with the same values on $S^n$. The forms $q^{D-e}H$ and
$\sum_iE_i^2+q\sum_iO_i^2$ of degree $2D$ agree on $S^n$, hence everywhere by
homogeneity:
\[
 q^{D-e}H=\sum_iE_i^2+\sum_i\sum_{j=0}^n(y_jO_i)^2 .
\]
Setting $y_0=1$ gives $q(1,X)=\omega(X)$ and $H(1,X)=f(X)$, so
$\omega^{D-e}f$ is a sum of squares in $\Q[X]$. Multiplying by further powers
of the rational sum of squares $\omega$ shows that every larger exponent also
works.

\subsection{Proof of Lemma~\ref{lem:fields-lift}}
\label{app:fields-lift}

Let $\mu=\mu_A/2$, so that $A\succeq2\mu I$, and $\nabla^2f\succeq2\mu I$. Put
$r_{ij}=y_iy_j-z_{ij}$ for $i\le j$. In each cubic monomial of the gradient
$\nabla f(y)$ replace the product $y_iy_j$ of its two smallest indices,
$i\le j$, by $z_{ij}$, and keep the terms of degree at most two. This gives a
rational quadratic map $a(y,z)$ with
$a(y,(y_iy_j)_{i\le j})=\nabla f(y)$ and
\[
 \nabla f(y)-a(y,z)=E(y)\,r(y,z),
\]
where the matrix $E(y)$ has homogeneous linear entries: a cubic term
$c\,y_ly_iy_j$ is replaced by $c\,y_lz_{ij}$, and the difference is
$c\,y_lr_{ij}$.

Put $d=x-y$ and let $\mathbf w=(d,\ y\otimes d,\ d\otimes d)$, of length
$n+2n^2$. Partitioning $A$ into its blocks $A_{00},A_{01},A_{11}$ for the
coordinates $v$ and $X\otimes v$, Lemma~\ref{lem:taylor-sos}(a) with the
polynomial center $y$, equivalently Taylor's formula with integral
remainder, gives
\[
 f(x)-f(y)-\nabla f(y)^{\mathsf T}d=\mathbf w^{\mathsf T}\mathcal M\mathbf w,\qquad
 \mathcal M=\begin{pmatrix}
 A_{00}/2&A_{01}/2&A_{01}/6\\
 A_{10}/2&A_{11}/2&A_{11}/6\\
 A_{10}/6&A_{11}/6&A_{11}/12\end{pmatrix},
\]
since $w(y+td,d)=(d,y\otimes d)+t(0,d\otimes d)$ and
$\int_0^1(1-t)t^j\,dt=\frac12,\frac16,\frac1{12}$ for $j=0,1,2$. The same
integration applied to $A\succeq2\mu I$ gives
$\mathcal M\succeq2\mu\bigl[\frac12I_n\oplus
\bigl(\left(\begin{smallmatrix}1/2&1/6\\1/6&1/12\end{smallmatrix}\right)\otimes
I_{n^2}\bigr)\bigr]$, and the $2\times2$ matrix has determinant $1/72>0$.
Hence $\mathcal N=\mathcal M-\frac\mu2\diag(I_n,0,0)\succ0$. Let $\mathcal D$ be
the rational matrix with $(\nabla f(y)-a)^{\mathsf T}d=r^{\mathsf T}\mathcal D\mathbf w$;
only the columns of $y\otimes d$ are used. Let
$\delta=\mu_{\mathcal N}$ and
$\lambda=\lceil1+\norm{\mathcal D}_F^2/(4\delta)\rceil$. The matrix
\[
 \mathcal K=\begin{pmatrix}\mathcal N&\mathcal D^{\mathsf T}/2\\
   \mathcal D/2&\lambda I\end{pmatrix}
\]
is positive definite, its Schur complement being
$\lambda I-\frac14\mathcal D\mathcal N^{-1}\mathcal D^{\mathsf T}\succeq I$. Put
\[
 g(y,z)=\frac{\norm{a(y,z)}^2}{2\mu}-f(y)+\lambda\norm{r(y,z)}^2 .
\]
We claim
\[
 f(x)+g(y,z)=(\mathbf w,r)^{\mathsf T}\mathcal K(\mathbf w,r)
   +\frac\mu2\Bigl\lVert d+\frac{a(y,z)}\mu\Bigr\rVert^2 .
\]
The first term equals $f(x)-f(y)-\nabla f(y)^{\mathsf T}d-\frac\mu2\norm d^2
+(\nabla f(y)-a)^{\mathsf T}d+\lambda\norm r^2$, and the second equals
$\frac\mu2\norm d^2+a^{\mathsf T}d+\norm a^2/(2\mu)$; their sum is
$f(x)-f(y)+\lambda\norm r^2+\norm a^2/(2\mu)=f(x)+g(y,z)$. Both matrices are
rational and positive semidefinite, so this is a rational sum of squares with
quadratic factors. Substituting $x=p$ gives $g\ge0$, because $f(p)=0$. At
$(y,z)=(p,(p_ip_j))$ we have $r=0$ and $a=\nabla f(p)=0$, so $g=0$. All
quantities are obtained by rational linear algebra of polynomial size from
$(f,A)$, without computing $p$.

\subsection{Proof of Theorem~\ref{thm:fields-blocks-height}}
\label{app:fields-blocks}

(a) The quartic $F_k^\circ$ is a sum of $2k+1$ squares of rational
quadratics with unique zero $p_k^\circ$, and $\delta_k=u_k(p_k^\circ)$ satisfies
$0<\delta_k\le M_k^{-2^k}$ (Section~\ref{sec:heights-witness}). With
$\lambda=\lceil3/\mu_{A_k^\circ}\rceil$, the rational matrix
$\lambda A_k^\circ+2ee^{\mathsf T}\succeq3I$ is a full Hessian Gram of $h_k$,
where $e$ selects the direction coordinate of the variable in $u_k$. The
quartic $h_k$ is a rational sum of squares, and it is positive: a zero would
be a zero of $F_k^\circ$, hence $p_k^\circ$, where $u_k^2=\delta_k^2>0$. Strong
convexity makes its minimum $m_k$ attained, and
$m_k\le h_k(p_k^\circ)=\delta_k^2\le M_k^{-2^{k+1}}$. The blocks $f$ and $\tilde g$
have minimum zero, attained at $p_1$ and $w_*$. The Hessian of $P_k$ is block
diagonal, so adding the three block Hessian certificates gives a rational
sum-of-squares Hessian certificate and $\nabla^2P_k\succeq I$; adding the
rational sums of squares of $f+\tilde g$ and $h_k$ gives a rational sum of
squares of $P_k$. The variable count is $3+9+2k$.

(b) Choose rationals $c$ and $s$ with $0<c$, $0<s$, and $c+s<m_k$, and a
rational point $x'$ with $f(x')<s$, possible since $f(p_1)=0$. By
Lemma~\ref{lem:taylor-sos}(c), applied at a rational center close to the
minimizer, the certified quartics $f+c$ and $h_k-c-s$, whose minima are
positive, have rational positive definite Gram matrices, hence rational sums of
squares. Specializing the joint rational certificate of $f+\tilde g$ at
$x=x'$ shows that $f(x')+\tilde g$ is a rational sum of squares, and so is
$\tilde g+s=(f(x')+\tilde g)+(s-f(x'))$. Then
$P_k=(f+c)+\bigl[(\tilde g+s)+(h_k-c-s)\bigr]$ is block separated. No size
bound is claimed for these certificates.

(c) Group a separated certificate as $P_k=S_x(x)+S'(w,\xi)$ with $S_x$ and
$S'$ nonnegative. As in the proof of Theorem~\ref{thm:fields-blocks},
$S_x=f+c$ and $S'=\tilde g+h_k-c$ for a rational constant $c$. Evaluating
$S_x$ at $p_1$ gives $c\ge0$, and evaluating $S'$ at $w_*$ and a minimizer of
$h_k$ gives $c\le m_k$. If $c=0$, then $S_x$ would be a rational certificate
of $f$, which Theorem~\ref{thm:fields-tower}(b) with $E=\Q$ excludes. Thus
$0<c\le M_k^{-2^{k+1}}$, and writing $c$ in lowest terms with a positive
numerator shows $\operatorname{den}(c)\ge M_k^{2^{k+1}}$. For a separated Gram
certificate, the constant entry of $Q_x$ equals $f(0)+c$, so
$\operatorname{den}(c)\le b_0\operatorname{den}((Q_x)_{00})$. For a separated
sum of squares with $x$-factors $u_j$, $c=\sum_ju_j(0)^2-f(0)$, so
$\operatorname{den}(c)\le b_0\prod_j\operatorname{den}(u_j(0))^2$. Taking
logarithms gives both bounds, and
$2^k\log_2M_k=(k+3)2^k\log_21000=\Omega(k2^k)$.
